%==============================================================================
%  CAPÍTULO — LOS MODELOS LATINOAMERICANOS
%==============================================================================
\chapter{Los modelos latinoamericanos}
\label{cap:comparado-latam}

\fichaCapitulo{El vecindario normativo de la insolvencia chilena.}{Argentina, Brasil, Colombia,
Perú, México y Uruguay: soluciones regionales al concurso de la empresa y de la persona natural,
y su contraste con la Ley \Nro 20.720.}

%==============================================================================
\bloqueTeorico
%==============================================================================

\subsection{Un vecindario con soluciones divergentes}

América Latina no tiene un derecho concursal uniforme. Pese a compartir tradición jurídica,
estructura económica y problemas semejantes ---informalidad, predominio de la micro y pequeña
empresa, sobreendeudamiento de los hogares---, los ordenamientos de la región han adoptado
soluciones marcadamente distintas en tres cuestiones decisivas: quién puede acceder al concurso, si
el procedimiento es judicial o administrativo, y si la persona natural tiene derecho a una descarga
de saldos.

Esta divergencia hace del comparado regional un ejercicio especialmente fértil para Chile. A
diferencia del contraste con Estados Unidos ---donde la distancia institucional obliga a traducir
categorías---, aquí se comparan sistemas que enfrentan las mismas restricciones presupuestarias y
la misma estructura empresarial.

\begin{foco}[La pregunta que ordena la comparación regional]
Toda la región enfrenta el mismo dilema: un procedimiento judicial garantista es caro y lento, y
por eso inalcanzable para la microempresa y para el deudor persona natural ---que son la inmensa
mayoría de los insolventes reales---. Colombia y Chile respondieron
\destacar{desjudicializando}; Perú, creando una \destacar{autoridad administrativa concursal};
Brasil y Argentina, manteniendo la vía judicial y simplificándola. Cada modelo paga un precio
distinto en garantías, costo y celeridad.
\end{foco}

%==============================================================================
\bloqueNormativo
%==============================================================================

\subsection{Argentina: la Ley \Nro 24.522 y el \emph{cramdown} local}

La Ley de Concursos y Quiebras argentina estructura el sistema en dos procedimientos clásicos: el
\destacar{concurso preventivo} y la \destacar{quiebra}. Tres rasgos merecen atención comparada:

\paragraph{El período de exclusividad.} El deudor dispone de un plazo para formular y obtener las
conformidades a su propuesta, agrupando a los acreedores en categorías. La aprobación requiere
mayoría de acreedores y de capital dentro de cada categoría, esquema muy próximo a la doble mayoría
del \art{79} chileno.

\paragraph{El \emph{salvataje} o \emph{cramdown} del artículo 48.} Es la institución argentina de
mayor originalidad. Fracasada la propuesta del deudor, en ciertas sociedades se abre un
procedimiento en que \destacar{terceros y acreedores pueden competir por adquirir la empresa},
formulando sus propias propuestas a los acreedores y desplazando a los socios originales si
obtienen las conformidades. Traslada la titularidad de la empresa a quien logre el acuerdo, en
lugar de liquidarla.

\paragraph{El privilegio laboral y el \emph{pronto pago}.} La ley contempla un mecanismo de
\destacar{pronto pago} de los créditos laborales, que permite satisfacerlos sin esperar el trámite
completo de verificación. Es el equivalente funcional del pago administrativo del \art{244}
chileno.

\begin{alerta}[El salvataje argentino y el vacío chileno]
La \LIR{} no contempla nada semejante al artículo 48 argentino. Fracasado el acuerdo, Chile va
derecho a la liquidación (\art{96}), aunque existan terceros dispuestos a capitalizar la empresa.
Sólo el rescate del \art{257} ---ya abierta la liquidación--- permite algo parecido, y siempre por
iniciativa del deudor. La posibilidad de que un tercero compita por la empresa \emph{antes} de
liquidarla es una vía de conservación que el legislador chileno no ha explorado.
\end{alerta}

\subsection{Brasil: recuperación judicial, extrajudicial y la reforma de 2020}

La Ley \Nro 11.101 de 2005 sustituyó el antiguo \emph{concordato} por un sistema de tres vías:
\destacar{recuperação judicial}, \destacar{recuperação extrajudicial} y \destacar{falência}. La
reforma de la Ley \Nro 14.112 de 2020 la modernizó profundamente.

\paragraph{La recuperación judicial.} El deudor obtiene un \emph{stay period} ---suspensión de
ejecuciones--- y presenta un plan que se vota por clases: laborales, garantizados, quirografarios y
microempresas. El plan requiere aprobación de todas las clases, con la particularidad de que la
clase laboral vota \destacar{sólo por cabeza}, sin ponderar el monto del crédito.

\paragraph{El \emph{cram down} brasileño.} Rechazado el plan por alguna clase, el juez puede
homologarlo si se cumplen requisitos cuantitativos tasados de apoyo global y de apoyo mínimo en la
clase disidente, más la ausencia de trato discriminatorio dentro de esa clase. Brasil, a diferencia
de Chile, \destacar{sí dispone de arrastre entre clases}.

\paragraph{La reforma de 2020.} Introdujo el financiamiento con prioridad (\emph{DIP financing}),
la mediación concursal, la insolvencia transfronteriza sobre la base de la Ley Modelo
\textsc{cnudmi} ---como el Capítulo VIII chileno--- y reglas sobre grupos empresariales, materia en
que el derecho chileno carece de regulación específica.

\begin{practica}[La consolidación de grupos: la laguna chilena]
Brasil reguló expresamente la consolidación procesal y sustancial de grupos empresariales; Chile no
tiene norma alguna al respecto, pese a que la insolvencia de holdings es cotidiana. Ante un grupo
en crisis, el abogado chileno debe articular procedimientos separados y coordinar de hecho lo que
la ley no coordina de derecho. Es una de las carencias más sentidas en la práctica de concursos de
envergadura.
\end{practica}

\subsection{Colombia: la Ley 1116 y la insolvencia de la persona natural}

Colombia ofrece, para Chile, el comparador regional más pertinente por dos razones.

\paragraph{El régimen empresarial (Ley 1116 de 2006).} Contempla la
\destacar{reorganización} y la \destacar{liquidación judicial}, con una particularidad
institucional relevante: la competencia se reparte entre los jueces civiles y la
\destacar{Superintendencia de Sociedades}, que ejerce funciones jurisdiccionales. Esa atribución de
jurisdicción a un órgano administrativo especializado no tiene equivalente chileno ---la \superir{}
fiscaliza y tramita la renegociación, pero no juzga---, y explica buena parte de la celeridad del
sistema colombiano.

\paragraph{La insolvencia de la persona natural no comerciante.} Regulada en el Código General del
Proceso, opera ante \destacar{centros de conciliación y notarios}, con una secuencia de negociación
de deudas, convalidación de acuerdo privado y, en su defecto, liquidación patrimonial. Es un
esquema de desjudicialización comparable al de la renegociación chilena, pero apoyado en
conciliadores privados en lugar de un órgano estatal.

\begin{foco}[Dos maneras de desjudicializar]
Chile y Colombia llegaron a la misma conclusión ---la persona natural no puede litigar su
insolvencia ante un juez--- y la resolvieron distinto. Chile creó un \destacar{procedimiento
administrativo gratuito} ante la \superir{}, financiado con presupuesto público. Colombia lo
entregó a \destacar{conciliadores y notarios}, con costo para el deudor pero sin carga fiscal. Los
datos chilenos de acceso masivo (capítulo \ref{cap:renegociacion}) sugieren que la gratuidad es un
determinante decisivo de la cobertura efectiva.
\end{foco}

\subsection{Perú: el modelo administrativo del \textsc{indecopi}}

Perú representa el extremo de la desjudicialización regional. El Sistema Concursal de la Ley
General del Sistema Concursal se tramita ante el \destacar{\textsc{indecopi}} ---el Instituto
Nacional de Defensa de la Competencia y de la Protección de la Propiedad Intelectual---, autoridad
administrativa que conoce del procedimiento concursal ordinario y del preventivo.

La \destacar{Junta de Acreedores} peruana concentra un poder decisorio muy superior al de su
homóloga chilena: decide el destino del deudor ---reestructuración o liquidación---, aprueba el
plan y designa al administrador o liquidador. El rol de la autoridad es de conducción del
procedimiento y control de legalidad, no de decisión sobre el fondo.

\begin{alerta}[Lo que Chile puede y no puede tomar del modelo peruano]
La atribución de todo el procedimiento a una autoridad administrativa reduce costos y tiempos, pero
plantea un problema constitucional que en Chile sería serio: la decisión sobre derechos
patrimoniales ---y sobre la extinción de créditos--- por un órgano que no ejerce jurisdicción. El
diseño chileno, que reserva a la \superir{} la renegociación y entrega al juez la liquidación y la
reorganización, es más conservador pero también más resistente al reproche de inconstitucionalidad
(capítulo \ref{cap:institucionalidad}).
\end{alerta}

\subsection{México y Uruguay: notas breves}

\paragraph{México.} La Ley de Concursos Mercantiles estructura un procedimiento en dos etapas
sucesivas ---\destacar{conciliación} y, fracasada ésta, \destacar{quiebra}---, con la
particularidad institucional del \destacar{\textsc{ifecom}}, órgano auxiliar del Poder Judicial
Federal encargado de designar y supervisar a los especialistas: visitador, conciliador y síndico.
Esa tripartición de auxiliares según la etapa contrasta con la dualidad chilena de veedor y
liquidador.

\paragraph{Uruguay.} La Ley \Nro 18.387 adoptó un \destacar{procedimiento único} de concurso, con
una solución de técnica legislativa notable: la calificación del concurso como
\destacar{culpable o fortuito}, con consecuencias sobre la responsabilidad de los administradores y
sobre los efectos para el deudor. Esa categoría ---inspirada en el derecho español--- cumple una
función análoga a la del incidente de mala fe chileno del \art{169 A}, pero opera como calificación
general del concurso y no como incidente sobre la descarga.

%==============================================================================
\bloquePractico
%==============================================================================

\subsection{Cuadro comparativo regional}

\begin{table}[htbp]
\centering
\footnotesize
\caption{Rasgos estructurales de los sistemas concursales latinoamericanos}
\label{tab:cap27-latam}
\begin{tabularx}{\textwidth}{@{}>{\raggedright\arraybackslash}p{0.13\textwidth}
                                 >{\raggedright\arraybackslash}p{0.19\textwidth}
                                 >{\raggedright\arraybackslash}p{0.17\textwidth} X@{}}
\toprule
\textbf{País} & \textbf{Órgano competente} & \textbf{Arrastre de clases} & \textbf{Persona natural} \\
\midrule
Chile & Juez civil; \superir{} en renegociación & No & Renegociación administrativa y
liquidación con \discharge{} \\
\addlinespace
Argentina & Juez comercial & Limitado; \emph{salvataje} del art.\ 48 & Sin régimen general de
descarga \\
\addlinespace
Brasil & Juez & Sí, con requisitos tasados & Régimen incorporado por la reforma de 2020 \\
\addlinespace
Colombia & Juez civil y Superintendencia de Sociedades & Restringido & Ante conciliadores y
notarios \\
\addlinespace
Perú & \textsc{indecopi} (administrativo) & Decide la junta & Comprendida en el sistema
concursal \\
\addlinespace
México & Juez de distrito, con \textsc{ifecom} & Limitado & Sin régimen desarrollado \\
\addlinespace
Uruguay & Juez concursal & Limitado & Concurso único con calificación \\
\bottomrule
\end{tabularx}
\end{table}

\verificar{Actualizar el cuadro contra el texto vigente de cada ordenamiento antes de publicar: la
región ha legislado intensamente en insolvencia desde 2020 y varias de estas reglas han sido
modificadas.}

\begin{foco}[La posición de Chile en la región]
El sistema chileno destaca en dos dimensiones y queda rezagado en una tercera. Destaca por la
\destacar{cobertura de la persona natural} ---la renegociación administrativa y gratuita es, en
acceso efectivo, de las más eficaces de la región--- y por la \destacar{institucionalidad
fiscalizadora} de la \superir{}. Queda rezagado en el instrumento de reestructuración empresarial:
sin arrastre entre clases, sin salvataje por terceros y sin regulación de grupos, la reorganización
chilena es comparativamente débil, lo que se refleja en las cifras ---apenas 32 reorganizaciones
en nueve meses frente a casi ocho mil procedimientos de persona deudora (capítulo
\ref{cap:renegociacion})---.
\end{foco}
